\SourceLeaf{060}
\begin{center}
\pdfbookmark[2]{42.}{article-042}42.
\end{center}
Wir haben die bisher beschriebene Tafel über einen geeigneten Bereich berechnet und diesem Werk beigefügt (Tafel I). Auf die Ellipse bezieht sich allein der erste Teil. Den zweiten Teil, der die hyperbolische Bewegung umfasst, erläutern wir weiter unten. Das Argument der Tafel, die Größe $A$, geht in Schritten von einem Tausendstel von $0$ bis $0{,}300$. Es folgen $\log B$ und $C$, deren Angaben in $10\,000\,000^{\mathrm{steln}}$, also mit sieben Dezimalstellen, zu lesen sind. Die ersten Ziffern vor den signifikanten Stellen sind nämlich weggelassen. Die vierte Spalte enthält schließlich die Größe $T$, zunächst auf 5, dann auf 6 Stellen berechnet. Diese Genauigkeit genügt vollkommen, da die Spalte nur dazu dient, die zum Argument $T$ gehörigen Werte von $\log B$ und $C$ zu finden, wenn man nach dem vorigen Artikel $t$ aus $v$ bestimmen will. Die umgekehrte Aufgabe, die sehr viel häufiger vorkommt, besteht darin, $v$ und $r$ aus $t$ zu bestimmen, und lässt sich ganz ohne die Größe $T$ lösen. Deshalb haben wir als Argument unserer Tafel lieber $A$ gewählt als $T$, das sich sonst beinahe ebenso gut geeignet und die Berechnung der Tafel sogar etwas erleichtert hätte. Es sei darauf hingewiesen, dass alle Zahlen der Tafel ursprünglich auf zehn Stellen berechnet wurden, sodass man sich überall auf die sieben hier angegebenen Stellen verlassen kann. Auf die bei dieser Arbeit verwendeten analytischen Verfahren können wir hier nicht näher eingehen, da uns eine ausführliche Erläuterung zu weit von unserem Vorhaben abbringen würde. Der Bereich der Tafel genügt vollkommen für alle Fälle, in denen das bisher erläuterte Verfahren vorteilhaft ist. Jenseits der Grenze $A=0{,}3$, der $T=0{,}392374$ oder $E=64^0\,8'$ entspricht, kann man, wie oben gezeigt, ohne Weiteres auf besondere Verfahren verzichten.
